\documentclass[a4paper,12pt]{article}
\usepackage[utf8]{inputenc}
\usepackage[T1]{fontenc}
\usepackage[magyar]{babel}
\usepackage{amsmath, amssymb, amsthm}
\usepackage{geometry}
\usepackage{xcolor}
\usepackage{tcolorbox}
\tcbuselibrary{theorems,skins,breakable}
\usepackage{tikz}
\usepackage{enumitem}
\usepackage{hyperref}
\geometry{margin=2.5cm}

% Egyedi stílusú dobozok a mikrotananyag módszertanhoz
\newtcolorbox{definiciodoboz}[1][]{colback=blue!5!white,colframe=blue!75!black,fonttitle=\bfseries,title=Definíció: #1}
\newtcolorbox{teteldoboz}[1][]{colback=green!8!white,colframe=green!75!black,fonttitle=\bfseries,title=Tétel / Szabály: #1}
\newtcolorbox{feladatdoboz}[1][]{colback=orange!5!white,colframe=orange!85!black,fonttitle=\bfseries,title=Ellenőrző feladat: #1}
\newtcolorbox{megoldasdoboz}[1][]{colback=gray!5!white,colframe=gray!60!black,fonttitle=\bfseries,title=Megoldás és magyarázat: #1}

\title{Témakör -- Mikrotananyag}
\author{AI LaTeX Studio}
\date{\today}

\begin{document}
\maketitle

\section{1. Mikro-egység: Alapfogalom bemutatása}

\begin{definiciodoboz}[Fogalom meghatározása]
Itt szerepel a témakör legkisebb, önállóan megérthető egysége, szabálya vagy definíciója. Egyszerre mindig csak egy új gondolatot vezetünk be.
\end{definiciodoboz}

% Azonnali visszacsatolás: Elmélet 1 -> Feladat 1
\begin{feladatdoboz}[1. Megértést ellenőrző kérdés (Feleletválasztós)]
Melyik állítás írja le helyesen a fenti fogalmat?
\begin{enumerate}[label=(\Alph*)]
    \item Első válaszlehetőség
    \item Második válaszlehetőség (Helyes megoldás)
    \item Harmadik válaszlehetőség
    \item Negyedik válaszlehetőség
\end{enumerate}
\end{feladatdoboz}

\begin{megoldasdoboz}[Megoldás az 1. kérdéshez]
\textbf{Helyes válasz: (B)} \\
\textit{Indoklás:} A definíció alapján ez az állítás tükrözi a pontos összefüggést.
\end{megoldasdoboz}

\section{2. Mikro-egység: Szabály és összefüggés}

\begin{teteldoboz}[Alapvető tétel / képlet]
A következő logikus lépésként bemutatott alaptétel vagy számítási képlet:
\[
    f(x) = a \cdot x + b
\]
ahol $a$ a meredekség és $b$ a tengelymetszet.
\end{teteldoboz}

% Azonnali visszacsatolás: Elmélet 2 -> Feladat 2
\begin{feladatdoboz}[2. Megértést ellenőrző feladat (Igaz/Hamis és Többlépcsős vezetés)]
\textbf{A) Igaz vagy Hamis?}
\begin{itemize}
    \item Ha $a > 0$, akkor a függvény szigorúan monoton növekvő. (I / H)
\end{itemize}

\medskip
\textbf{B) Többlépcsős feladatvezetés:}
\begin{enumerate}[label=\alph*)]
    \item Számítsd ki a függvényértéket $x = 3$ esetén, ha $a = 2$ és $b = 4$!
    \item Határozd meg, melyik pontban metszi az $y$-tengelyt!
\end{enumerate}
\end{feladatdoboz}

\begin{megoldasdoboz}[Megoldás a 2. feladathoz]
\textbf{A) Igaz}, mert pozitív meredekség esetén az argumentum növekedésével a függvényérték is nő. \\
\textbf{B) Lépésről lépésre:}
\begin{enumerate}[label=\alph*)]
    \item $f(3) = 2 \cdot 3 + 4 = 6 + 4 = 10$.
    \item $x = 0$ helyettesítéssel $f(0) = 4$, tehát a metszéspont $(0, 4)$.
\end{enumerate}
\end{megoldasdoboz}

\end{document}
